% GENERATED FILE. Edit canonical lessons, then run npm run generate:books.
% canonical-source-hash: fnv1a64:ec5489b1e125c074
% canonical-lessons: IT-C286-edizione, IT-C286-didascalia, IT-C286-scoop, IT-C286-censura, IT-C286-censurare

\chapter{Carta stampata}
\label{ch:it-carta-stampata}
\hlchaptermodality{286}

\begin{chapteropening}
\textbf{By the end of this chapter:} \emph{I can talk about what gets printed: the edition, the caption, the scoop and censorship.}

Complete the last lesson of chapter 286: I can talk about what gets printed: the edition, the caption, the scoop and censorship.
\end{chapteropening}

\section[l'edizione]{l'edizione — the edition; the news bulletin}

\label{lesson:IT-C286-edizione}



\begin{quote}
Before the new one: say the Italian for the press release, then the Italian for the denial.
\end{quote}

\subsection*{You'll want to know: l'edizione}
\textbf{l'edizione} — \textquotedblleft{}the edition; the news bulletin\textquotedblright{}.

\begin{quote}
L'edizione della sera ha una foto in prima pagina.
\end{quote}

\begin{quote}
\emph{The evening edition has a photo on the front page.}
\end{quote}

\emph{Un'edizione straordinaria} — a special edition. On TV, \emph{l'edizione delle venti} is the eight o'clock news.

\begin{cousinweb}
\emph{edizione} is a learned borrowing from Latin \emph{ēditiō} \textquotedblleft{}a putting forth, a publishing\textquotedblright{} (Lewis \& Short), from \emph{ēdere} \textquotedblleft{}to give out\textquotedblright{}.
\end{cousinweb}

\begin{grammarlens}[title={What you've built}]
One more word for telling, explaining and describing at B1.
\end{grammarlens}

\subsection*{Your turn}
\begin{itemize}
\raggedright
  \item \emph{Say it:} \emph{l'edizione}
  \item \emph{Say it:} \emph{l'edizione}, once more
  \item \emph{Say it:} a sentence of your own with \emph{l'edizione}
  \item \emph{Recall:} say the Italian for the press release, then the Italian for the denial, then say \emph{l'edizione} again
\end{itemize}

\begin{tcolorbox}[breakable,colback=teal!4,colframe=teal!35!black,title={Before you move on}]
What does l'edizione mean? (\textquotedblleft{}The edition; the news bulletin\textquotedblright{}.) Say it once more.
\end{tcolorbox}
\section[la didascalia]{la didascalia — the caption}

\label{lesson:IT-C286-didascalia}



\begin{quote}
Before the new one: say the Italian for the denial, then the Italian for the edition.
\end{quote}

\subsection*{You'll want to know: la didascalia}
\textbf{la didascalia} — \textquotedblleft{}the caption\textquotedblright{}.

\begin{quote}
La didascalia sotto la foto è sbagliata.
\end{quote}

\begin{quote}
\emph{The caption under the photo is wrong.}
\end{quote}

Also the stage directions in a play.

\begin{cousinweb}
\emph{didascalia} is a learned word usually traced to Greek \emph{didaskalía} \textquotedblleft{}teaching\textquotedblright{}, built on \emph{didáskein} \textquotedblleft{}to teach\textquotedblright{}; EtymDB has no Italian entry.
\end{cousinweb}

\begin{grammarlens}[title={What you've built}]
One more word for telling, explaining and describing at B1.
\end{grammarlens}

\subsection*{Your turn}
\begin{itemize}
\raggedright
  \item \emph{Say it:} \emph{la didascalia}
  \item \emph{Say it:} \emph{la didascalia}, once more
  \item \emph{Say it:} a sentence of your own with \emph{la didascalia}
  \item \emph{Recall:} say the Italian for the denial, then the Italian for the edition, then say \emph{la didascalia} again
\end{itemize}

\begin{tcolorbox}[breakable,colback=teal!4,colframe=teal!35!black,title={Before you move on}]
What does la didascalia mean? (\textquotedblleft{}The caption\textquotedblright{}.) Say it once more.
\end{tcolorbox}
\section[lo scoop]{lo scoop — the scoop, the exclusive story}

\label{lesson:IT-C286-scoop}



\begin{quote}
Before the new one: say the Italian for the edition, then the Italian for the caption.
\end{quote}

\subsection*{You'll want to know: lo scoop}
\textbf{lo scoop} — \textquotedblleft{}the scoop, the exclusive story\textquotedblright{}.

\begin{quote}
Quel giornalista ha fatto uno scoop.
\end{quote}

\begin{quote}
\emph{That journalist got a scoop.}
\end{quote}

Borrowed from English; plural unchanged: \emph{gli scoop}.

\begin{cousinweb}
\emph{scoop} was borrowed from English \emph{scoop} (EtymDB), a shovel; in the press, a story scooped up before anyone else.
\end{cousinweb}

\begin{grammarlens}[title={What you've built}]
One more word for telling, explaining and describing at B1.
\end{grammarlens}

\subsection*{Your turn}
\begin{itemize}
\raggedright
  \item \emph{Say it:} \emph{lo scoop}
  \item \emph{Say it:} \emph{lo scoop}, once more
  \item \emph{Say it:} a sentence of your own with \emph{lo scoop}
  \item \emph{Recall:} say the Italian for the edition, then the Italian for the caption, then say \emph{lo scoop} again
\end{itemize}

\begin{tcolorbox}[breakable,colback=teal!4,colframe=teal!35!black,title={Before you move on}]
What does lo scoop mean? (\textquotedblleft{}The scoop, the exclusive story\textquotedblright{}.) Say it once more.
\end{tcolorbox}
\section[la censura]{la censura — censorship}

\label{lesson:IT-C286-censura}



\begin{quote}
Before the new one: say the Italian for the caption, then the Italian for the scoop.
\end{quote}

\subsection*{You'll want to know: la censura}
\textbf{la censura} — \textquotedblleft{}censorship\textquotedblright{}.

\begin{quote}
Il film ha avuto problemi con la censura.
\end{quote}

\begin{quote}
\emph{The film had trouble with the censors.}
\end{quote}

Also a formal rebuke: \emph{un voto di censura}.

\begin{cousinweb}
\emph{censura} is a learned borrowing from Latin \emph{cēnsūra} \textquotedblleft{}the office of the censor\textquotedblright{} (Lewis \& Short), from \emph{cēnsēre} \textquotedblleft{}to assess, to judge\textquotedblright{}.
\end{cousinweb}

\begin{grammarlens}[title={What you've built}]
One more word for telling, explaining and describing at B1.
\end{grammarlens}

\subsection*{Your turn}
\begin{itemize}
\raggedright
  \item \emph{Say it:} \emph{la censura}
  \item \emph{Say it:} \emph{la censura}, once more
  \item \emph{Say it:} a sentence of your own with \emph{la censura}
  \item \emph{Recall:} say the Italian for the caption, then the Italian for the scoop, then say \emph{la censura} again
\end{itemize}

\begin{tcolorbox}[breakable,colback=teal!4,colframe=teal!35!black,title={Before you move on}]
What does la censura mean? (\textquotedblleft{}Censorship\textquotedblright{}.) Say it once more.
\end{tcolorbox}
\section[censurare]{censurare — to censor}

\label{lesson:IT-C286-censurare}



\begin{quote}
Before the new one: say the Italian for the scoop, then the Italian for censorship.
\end{quote}

\subsection*{You'll want to know: censurare}
\textbf{censurare} — \textquotedblleft{}to censor\textquotedblright{}.

\begin{quote}
Hanno censurato una scena del film.
\end{quote}

\begin{quote}
\emph{They censored a scene from the film.}
\end{quote}

In formal Italian it also means to criticise strongly.

\begin{cousinweb}
\emph{censurare} is built on \emph{censura} \textquotedblleft{}censorship\textquotedblright{}, from Latin \emph{cēnsūra} and \emph{cēnsēre} \textquotedblleft{}to assess, to judge\textquotedblright{} (Lewis \& Short).
\end{cousinweb}

\begin{grammarlens}[title={What you've built}]
One more word for telling, explaining and describing at B1.
\end{grammarlens}

\subsection*{Your turn}
\begin{itemize}
\raggedright
  \item \emph{Say it:} \emph{censurare}
  \item \emph{Say it:} \emph{censurare}, once more
  \item \emph{Say it:} a sentence of your own with \emph{censurare}
  \item \emph{Recall:} say the Italian for the scoop, then the Italian for censorship, then say \emph{censurare} again
\end{itemize}

\begin{tcolorbox}[breakable,colback=teal!4,colframe=teal!35!black,title={Before you move on}]
What does censurare mean? (\textquotedblleft{}To censor\textquotedblright{}.) Say it once more.
\end{tcolorbox}
